\vfill%
\pagebreak
\section{Conditional loop}%
\label{sec:flow:while_loop}

\token{while} evaluates its block repeatedly, as long as a condition of type
\token{bool} holds. Its variant \token{while let} repeats as long as a value fits
a pattern (Section~\ref{sec:flow:while_let_loop}).

\subsection{Grammar}

\begin{lstlisting}[style=grammarVerb]
while_expression := 'while' condition block
                  | 'while' 'let' pattern '=' expression guard? block
guard            := '&&' expression
\end{lstlisting}

As for \token{loop}, the body is always a block.

\subsection{Evaluation}

The condition is evaluated before each iteration, including the first. If it
holds, the block is evaluated, then the condition again; if it does not, the loop
ends. The block may therefore be evaluated zero times.

\token{break} and \token{continue} behave as in a \token{loop}
(Section~\ref{sec:flow:inf_loop}). After a \token{continue}, the condition is
evaluated again before the next iteration. The control flow of the next listing
is drawn in Figure~\ref{fig:flow:simple_while_loop}.

\begin{lstlisting}[style=coloredverbatim, caption=Simple while loop, label=lst:flow:simple_while_loop]
let mut count = 0;

while count < 5 {
  count += 1;
}

assert(count == 5);
\end{lstlisting}

\begin{figure}[h]
  \centering
  \begin{adjustbox}{max width=\linewidth}
    \begin{tikzpicture}[node distance=6mm]
      \node[cfterm] (start) {start of foo};
      \node[cfstep, below=of start] (init) {let mut count = 0;};
      \node[cftest, below=of init] (test) {count < 5};
      \node[cfstep, below=8mm of test] (body) {count += 1;};
      \node[cfstep, below=of body] (after) {assert(count == 5);};
      \node[cfterm, below=of after] (end) {end of foo};

      \draw[cfflow] (start) -- (init);
      \draw[cfflow] (init) -- (test);
      \draw[cfflow] (test) -- node[cflab, right] {true} (body);
      \draw[cfflow] (body.east) -- ++(14mm, 0) |- (test.east);
      \draw[cfflow] (test.west) -- node[cflab, above] {false} ++(-14mm, 0) |- (after.west);
      \draw[cfflow] (after) -- (end);
    \end{tikzpicture}
  \end{adjustbox}
  \caption{\label{fig:flow:simple_while_loop} The flow graph of the
    \token{while} loop in Listing~\ref{lst:flow:simple_while_loop}. The
    condition is evaluated before every iteration, including the first.}
\end{figure}


The condition must be of type \token{bool}, and must not be known at compile
time. A condition known to be \token{false} makes the loop dead code (E4143). A
condition known to be \token{true} makes the \token{while} an infinite loop,
which must be written \token{loop} (E4100).

\begin{lstlisting}[style=coloredverbatimError, escapechar=@]
let a = [1, 2, 3];
while @\hb{a.len > 5}@ { // error, loop is never entered
  println(a);
}

let mut i = 0;
while @\hb{true}@ { // error, use /loop/
  i += 1;
  if i == 10 { break; }
}
println(i);
\end{lstlisting}

\subsection{Loops entered at least once}
\label{sec:flow:loop_at_least_once}

Ymir has no loop that tests its condition after its body. A loop whose body must
run at least once is a \token{loop} (Section~\ref{sec:flow:inf_loop}) that breaks
at the end of its body.

\begin{lstlisting}[style=coloredverbatim, caption=A loop entered at least once, label=lst:flow:loop_at_least_once]
let mut count = 0;

loop {
  count += 1;
  if count >= 5 { break; }
}

assert(count == 5);
\end{lstlisting}

\subsection{Value and type}
\label{sec:flow:while_value}

A \token{while} has no value, and its type is \token{void}: the loop can end
because its condition failed, and no \token{break} then provides a value. A
\token{break} in a \token{while} carries no value; one that does is an error,
E4095. Only \token{loop} produces a value.

\subsection{While-let loop}
\label{sec:flow:while_let_loop}

\token{while let pattern = value} evaluates \token{value} before each iteration,
and evaluates the block if it fits \token{pattern}, with the variables of the
pattern in scope. The loop ends at the first evaluation that does not fit. A
\token{guard} is evaluated only when the pattern fits, and the loop ends if it
does not hold. In the following example, the loop continues as long as \token{x}
holds a value lower than \token{5}.

\begin{lstlisting}[style=coloredverbatim, caption=While-let example, label=lst:flow:while_let_example]
fn foo(i: i32)-> i32?;

fn bar() {
  let mut x = foo(0);
  let mut count = 0;

  while let Ok(z) = x && z < 5 {
    count += 1;
    x = foo(count);
  }
  println(count);
}
\end{lstlisting}

The compiler rewrites a \token{while let} into a \token{loop} around an
\token{if let} (Section~\ref{sec:flow:cond_pattern}) whose \token{else} breaks,
so Listing~\ref{lst:flow:while_let_example} and
Listing~\ref{lst:flow:while_let_rewritten} are the same program. Everything that
Section~\ref{sec:flow:cond_pattern} says of the pattern and the guard applies to
\token{while let}, including the requirement that the pattern be refutable.

\begin{lstlisting}[style=coloredverbatim, caption=While-let rewritten using loop and if-let, label=lst:flow:while_let_rewritten]
fn foo(i: i32)-> i32?;

fn bar() {
  let mut x = foo(0);
  let mut count = 0;

  loop {
    if let Ok(z) = x && z < 5 {
      count += 1;
      x = foo(count);
    } else break;
  }
  println(count);
}
\end{lstlisting}

\subsection{Diagnostics}

\begin{center}
  \begin{adjustbox}{max width=\linewidth}
    \begin{tabular}{|l|l|}
      \hline
      Code & Raised when\\[0pt]
      \hline
      \hline
      E4095 & the condition is not of type \token{bool}, or a \token{break} carries a value\\[0pt]
      E4100 & the condition is known to hold, or the pattern of a \token{while let} to fit\\[0pt]
      E4143 & the condition is known not to hold\\[0pt]
      \hline
    \end{tabular}
  \end{adjustbox}
\end{center}
